\nonstopmode{}
\documentclass[a4paper]{book}
\usepackage[times,inconsolata,hyper]{Rd}
\usepackage{makeidx}
\makeatletter\@ifl@t@r\fmtversion{2018/04/01}{}{\usepackage[utf8]{inputenc}}\makeatother
% \usepackage{graphicx} % @USE GRAPHICX@
\makeindex{}
\begin{document}
\chapter*{}
\begin{center}
{\textbf{\huge Package `ResIN'}}
\par\bigskip{\large \today}
\end{center}
\ifthenelse{\boolean{Rd@use@hyper}}{\hypersetup{pdftitle = {ResIN: Conduct Response Item Network (ResIN) analysis with social response data}}}{}
\ifthenelse{\boolean{Rd@use@hyper}}{\hypersetup{pdfauthor = {Philip Warncke; Dino Carpentras; Adrian Lüders}}}{}
\begin{description}
\raggedright{}
\item[Type]\AsIs{Package}
\item[Title]\AsIs{Conduct Response Item Network (ResIN) analysis with social
response data}
\item[Version]\AsIs{2.3.0}
\item[Author]\AsIs{Philip Warncke [aut, cre], Dino Carpentras [aut], Adrian Lüders [aut]}
\item[Maintainer]\AsIs{Philip Warncke }\email{philip.warncke@ul.ie}\AsIs{Maintainer: Philip Warncke }\email{philip.warncke@ul.ie}\AsIs{}
\item[Description]\AsIs{Contains various tools to perform and visualize Response Item Networks ('ResIN's'). 'ResIN' dummy-codes ordered and qualitative response choices from (survey) data, calculates pairwise associations and maps the location of each item response as a node in a force-directed network. Please refer to <}\url{https://www.resinmethod.net/}\AsIs{> for more details.}
\item[License]\AsIs{GPL-3}
\item[URL]\AsIs{}\url{https://pwarncke77.github.io/ResIN/}\AsIs{}
\item[BugReports]\AsIs{}\url{https://github.com/pwarncke77/ResIN/issues}\AsIs{}
\item[Depends]\AsIs{R (>= 4.2.0)}
\item[Imports]\AsIs{utils, graphics, psych, ggplot2 (>= 3.4.4), dplyr (>= 1.1.3),
tidyr (>= 1.3.0), fastDummies, qgraph, igraph,
DirectedClustering, foreach, parallelly, parallel, doSNOW,
readr, ggraph (>= 2.2.0), shadowtext (>= 0.1.4), tidygraph,
network, readr}
\item[Encoding]\AsIs{UTF-8}
\item[LazyData]\AsIs{true}
\item[LazyDataCompression]\AsIs{xz}
\item[RoxygenNote]\AsIs{7.3.3}
\item[Suggests]\AsIs{knitr, rmarkdown}
\item[VignetteBuilder]\AsIs{knitr}
\end{description}
\Rdcontents{Contents}
\HeaderA{as.gephi.ResIN}{Coerce a ResIN object to Gephi CSV table(s)}{as.gephi.ResIN}
%
\begin{Description}
Produces Gephi-readable edge (and optionally node) tables from a \code{ResIN}
object and (by default) writes them to CSV. Set \code{dont\_save\_csv = TRUE}
to return tables without writing files.
\end{Description}
%
\begin{Usage}
\begin{verbatim}
## S3 method for class 'ResIN'
as.gephi(
  x,
  file = "ResIN_gephi.csv",
  edges_only = TRUE,
  dont_save_csv = FALSE,
  weight_col = "weight",
  ...
)
\end{verbatim}
\end{Usage}
%
\begin{Arguments}
\begin{ldescription}
\item[\code{x}] A \code{ResIN} object.

\item[\code{file}] Output file name (legacy style).

\item[\code{edges\_only}] Logical; if TRUE write/return only edges.

\item[\code{dont\_save\_csv}] Logical; set TRUE to disable writing.

\item[\code{weight\_col}] Name of the edge-weight column in \code{x\$ResIN\_edgelist}.

\item[\code{...}] Ignored.
\end{ldescription}
\end{Arguments}
%
\begin{Value}
If \code{edges\_only = TRUE}, an edge table data.frame. Otherwise a list with \code{edges} and \code{nodes}. “When \code{dont\_save\_csv = FALSE}, the return value is returned invisibly.”
\end{Value}
%
\begin{Examples}
\begin{ExampleCode}
## Load the 12-item simulated Likert-type ResIN toy dataset
data(lik_data)

## Estimate a ResIN network
res <- ResIN(lik_data, plot_ggplot = FALSE)

## Create Gephi edge table without writing files
edges <- as.gephi(res, dont_save_csv = TRUE)
head(edges)

## Not run: 
## Write CSV file(s) for import to Gephi
## (writes "ResIN_gephi.csv" by default)
as.gephi(res, file = "ResIN_gephi.csv")

## Write both edges and nodes tables
## (writes "ResIN_gephi_edges.csv" and "ResIN_gephi_nodes.csv")
as.gephi(res, file = "ResIN_gephi.csv", edges_only = FALSE)

## End(Not run)


\end{ExampleCode}
\end{Examples}
\HeaderA{as.graphsjl}{Julia Graphs.jl coercion helpers}{as.graphsjl}
%
\begin{Description}
Generic for exporting objects to a lightweight Graphs.jl-ready representation.
For \code{ResIN} objects, this creates edge and node CSV tables suitable for
loading with \pkg{CSV.jl}/\pkg{DataFrames.jl}, with integer vertex IDs for
direct use in \pkg{Graphs.jl}.
\end{Description}
%
\begin{Usage}
\begin{verbatim}
as.graphsjl(x, ...)
\end{verbatim}
\end{Usage}
%
\begin{Arguments}
\begin{ldescription}
\item[\code{x}] Object to coerce/export.

\item[\code{...}] Passed to methods.
\end{ldescription}
\end{Arguments}
\HeaderA{as.graphsjl.ResIN}{Export a ResIN object to Graphs.jl (Julia) tables}{as.graphsjl.ResIN}
%
\begin{Description}
Produces Graphs.jl-style edge (and optionally node) tables from a \code{ResIN}
object. Edge endpoints are mapped to integer vertex IDs (\code{src}, \code{dst})
to align with \pkg{Graphs.jl} in Julia. Node and edge metadata are preserved as table
columns. The node table stores the vertex mapping and additional ResIN metadata.
\end{Description}
%
\begin{Usage}
\begin{verbatim}
## S3 method for class 'ResIN'
as.graphsjl(
  x,
  file = "ResIN_graphsjl.csv",
  edges_only = TRUE,
  dont_save_csv = FALSE,
  weight_col = "weight",
  ...
)
\end{verbatim}
\end{Usage}
%
\begin{Arguments}
\begin{ldescription}
\item[\code{x}] A \code{ResIN} object.

\item[\code{file}] Output file name (legacy style). If \code{edges\_only = TRUE},
the edge table is written to \code{file}. If \code{edges\_only = FALSE},
\code{file} is treated as a prefix and \code{"\_edges.csv"} /
\code{"\_nodes.csv"} are appended (with any trailing \code{.csv} removed).

\item[\code{edges\_only}] Logical; if TRUE (default), only write/return edge table.

\item[\code{dont\_save\_csv}] Logical; if FALSE (default), write CSV output. If TRUE,
no files are written and the resulting table(s) are returned visibly.

\item[\code{weight\_col}] Preferred edge-weight column name. Defaults to \code{"weight"}.

\item[\code{...}] Ignored.
\end{ldescription}
\end{Arguments}
%
\begin{Value}
If \code{edges\_only = TRUE}, an edge table \code{data.frame}.
Otherwise a list with elements \code{edges} and \code{nodes}. The node table
includes integer \code{vertex\_id} values and preserved node metadata.
\end{Value}
%
\begin{Examples}
\begin{ExampleCode}
## Not run: 
data(lik_data)
res <- ResIN(lik_data, generate_ggplot = FALSE, plot_ggplot = FALSE)

# Return tables only (no files written)
jl_tbls <- as.graphsjl(res, dont_save_csv = TRUE, edges_only = FALSE)

# Default behavior writes CSV files
# as.graphsjl(res, file = "ResIN_graphsjl.csv", edges_only = FALSE)

## End(Not run)

\end{ExampleCode}
\end{Examples}
\HeaderA{as.igraph.ResIN}{Coerce a ResIN object to an igraph graph}{as.igraph.ResIN}
%
\begin{Description}
Converts a \code{ResIN} object to an \code{igraph} graph using the adjacency
matrix stored in \code{x\$aux\_objects\$adj\_matrix}.
\end{Description}
%
\begin{Usage}
\begin{verbatim}
## S3 method for class 'ResIN'
as.igraph(x, mode = "undirected", weighted = TRUE, diag = FALSE, ...)
\end{verbatim}
\end{Usage}
%
\begin{Arguments}
\begin{ldescription}
\item[\code{x}] A \code{ResIN} object.

\item[\code{mode}, \code{weighted}, \code{diag}] Passed to \code{igraph::graph\_from\_adjacency\_matrix()}.

\item[\code{...}] Additional arguments passed to \code{igraph::graph\_from\_adjacency\_matrix()}.
\end{ldescription}
\end{Arguments}
%
\begin{Value}
An \code{igraph} object.
\end{Value}
%
\begin{Examples}
\begin{ExampleCode}
## Load the 12-item simulated Likert-type ResIN toy dataset
data(lik_data)

## Run the function:

igraph_output <-  as.igraph(ResIN(lik_data, plot_ggplot = FALSE))

class(igraph_output)

## Plot and/or investigate as you wish:

igraph::plot.igraph(igraph_output)


\end{ExampleCode}
\end{Examples}
\HeaderA{as.network.ResIN}{Convert a ResIN object to a statnet/network object}{as.network.ResIN}
%
\begin{Description}
Coerces a \code{ResIN} object to a \code{network} object (from the
\pkg{network} package used in the \pkg{statnet} ecosystem). The method
preserves edge-level columns from \code{x\$ResIN\_edgelist} as edge attributes
and node-level columns from \code{x\$ResIN\_nodeframe} as vertex attributes
whenever available.
\end{Description}
%
\begin{Usage}
\begin{verbatim}
## S3 method for class 'ResIN'
as.network(x, directed = FALSE, loops = FALSE, multiple = FALSE, ...)
\end{verbatim}
\end{Usage}
%
\begin{Arguments}
\begin{ldescription}
\item[\code{x}] A \code{ResIN} object.

\item[\code{directed}] Logical; should the resulting network be directed?
Defaults to \code{FALSE}.

\item[\code{loops}] Logical; allow self-loops? Defaults to \code{FALSE}.

\item[\code{multiple}] Logical; allow multiple edges? Defaults to \code{FALSE}.

\item[\code{...}] Additional arguments passed to \code{network::as.network()}.
\end{ldescription}
\end{Arguments}
%
\begin{Value}
An object of class \code{network}.
\end{Value}
%
\begin{Examples}
\begin{ExampleCode}
data(lik_data)
res <- ResIN(lik_data, generate_ggplot = FALSE, plot_ggplot = FALSE)

# ResIN re-exports network::as.network()
net <- as.network.ResIN(res) ## alternatively: as.network(res)
net

\end{ExampleCode}
\end{Examples}
\HeaderA{as.networkx}{Python NetworkX coercion helpers}{as.networkx}
%
\begin{Description}
Generic for exporting objects to a lightweight NetworkX-ready representation.
For \code{ResIN} objects, this creates edge and node CSV tables that can be
read via \pkg{pandas} and converted to a \pkg{networkx} graph.
\end{Description}
%
\begin{Usage}
\begin{verbatim}
as.networkx(x, ...)
\end{verbatim}
\end{Usage}
%
\begin{Arguments}
\begin{ldescription}
\item[\code{x}] Object to coerce/export.

\item[\code{...}] Passed to methods.
\end{ldescription}
\end{Arguments}
\HeaderA{as.networkx.ResIN}{Export a ResIN object to NetworkX (Python) tables}{as.networkx.ResIN}
%
\begin{Description}
Produces NetworkX-ready edge (and optionally node) tables
from a \code{ResIN} object for further manipulation in Python.
By default, this method writes CSV files that can be imported into Python
via \pkg{pandas} and \pkg{networkx}. Node and edge metadata are preserved as
table columns.
\end{Description}
%
\begin{Usage}
\begin{verbatim}
## S3 method for class 'ResIN'
as.networkx(
  x,
  file = "ResIN_networkx.csv",
  edges_only = TRUE,
  dont_save_csv = FALSE,
  weight_col = "weight",
  ...
)
\end{verbatim}
\end{Usage}
%
\begin{Arguments}
\begin{ldescription}
\item[\code{x}] A \code{ResIN} object.

\item[\code{file}] Output file name (legacy style). If \code{edges\_only = TRUE},
the edge table is written to \code{file}. If \code{edges\_only = FALSE},
\code{file} is treated as a prefix and \code{"\_edges.csv"} /
\code{"\_nodes.csv"} are appended (with any trailing \code{.csv} removed).

\item[\code{edges\_only}] Logical; if TRUE (default), only write/return edge table.

\item[\code{dont\_save\_csv}] Logical; if FALSE (default), write CSV output. If TRUE,
no files are written and the resulting table(s) are returned visibly.

\item[\code{weight\_col}] Preferred edge-weight column name. Defaults to \code{"weight"}.

\item[\code{...}] Ignored.
\end{ldescription}
\end{Arguments}
%
\begin{Value}
If \code{edges\_only = TRUE}, an edge table \code{data.frame}.
Otherwise a list with elements \code{edges} and \code{nodes}.
\end{Value}
%
\begin{Examples}
\begin{ExampleCode}
## Not run: 
data(lik_data)
res <- ResIN(lik_data, generate_ggplot = FALSE, plot_ggplot = FALSE)

# Return tables only (no files written)
nx_tbls <- as.networkx(res, dont_save_csv = TRUE, edges_only = FALSE)

# Default behavior writes CSV files
# as.networkx(res, file = "ResIN_networkx.csv", edges_only = FALSE)

## End(Not run)

\end{ExampleCode}
\end{Examples}
\HeaderA{as.qgraph.ResIN}{Coerce a ResIN object to a qgraph object}{as.qgraph.ResIN}
%
\begin{Description}
Converts a \code{ResIN} object to a \code{qgraph} object using the adjacency
matrix stored in \code{x\$aux\_objects\$adj\_matrix}.
\end{Description}
%
\begin{Usage}
\begin{verbatim}
## S3 method for class 'ResIN'
as.qgraph(
  x,
  layout = "spring",
  maximum = 1,
  vsize = 6,
  DoNotPlot = TRUE,
  sampleSize = NULL,
  title = "ResIN graph in qgraph",
  mar = c(3, 8, 3, 8),
  normalize = FALSE,
  ...
)
\end{verbatim}
\end{Usage}
%
\begin{Arguments}
\begin{ldescription}
\item[\code{x}] A \code{ResIN} object.

\item[\code{layout}, \code{maximum}, \code{vsize}, \code{DoNotPlot}, \code{sampleSize}, \code{title}, \code{mar}, \code{normalize}] Passed to \code{qgraph::qgraph()}.

\item[\code{...}] Additional arguments passed to \code{qgraph::qgraph()}.
\end{ldescription}
\end{Arguments}
%
\begin{Value}
A \code{qgraph} object.
\end{Value}
%
\begin{Examples}
\begin{ExampleCode}
## Load the 12-item simulated Likert-type ResIN toy dataset
data(lik_data)

## Run the function:
ResIN_qgraph <-  as.qgraph(ResIN(lik_data, plot_ggplot = FALSE))

class(ResIN_qgraph)

\end{ExampleCode}
\end{Examples}
\HeaderA{as.tidygraph.ResIN}{Coerce a ResIN object to a tidygraph graph}{as.tidygraph.ResIN}
%
\begin{Description}
Converts a \code{ResIN} object to a \code{tidygraph::tbl\_graph} object while
preserving as much node- and edge-level information as possible from
\code{x\$ResIN\_nodeframe} and \code{x\$ResIN\_edgelist}.

Because \code{tidygraph} stores edge endpoints as integer node indices, the
original edge endpoint labels are preserved in additional edge columns
\code{from\_name} and \code{to\_name}.

If \code{ResIN\_nodeframe} or \code{ResIN\_edgelist} are unavailable, the method
falls back to a simpler conversion via \code{as.igraph()} followed by
\code{tidygraph::as\_tbl\_graph()}, which may not preserve all metadata.
\end{Description}
%
\begin{Usage}
\begin{verbatim}
## S3 method for class 'ResIN'
as.tidygraph(x, directed = FALSE, ...)
\end{verbatim}
\end{Usage}
%
\begin{Arguments}
\begin{ldescription}
\item[\code{x}] A \code{ResIN} object.

\item[\code{directed}] Logical; should the resulting graph be treated as directed?
Defaults to \code{FALSE}.

\item[\code{...}] Ignored.
\end{ldescription}
\end{Arguments}
%
\begin{Value}
A \code{tidygraph::tbl\_graph} object. Node data include (when present)
all columns from \code{x\$ResIN\_nodeframe}; edge data include (when present)
all columns from \code{x\$ResIN\_edgelist}, plus \code{from\_name}/\code{to\_name}
preserving original endpoint labels.
\end{Value}
%
\begin{Examples}
\begin{ExampleCode}
## Load toy data and estimate ResIN
data(lik_data)
res <- ResIN(lik_data, network_stats = TRUE, detect_clusters = TRUE,
             plot_ggplot = FALSE)

## Convert to tidygraph
tg <- as.tidygraph(res)
tg

\end{ExampleCode}
\end{Examples}
\HeaderA{Bootstrap\_example}{Output example for a bootstrapping analysis conducted with the ResIN package}{Bootstrap.Rul.example}
\keyword{datasets}{Bootstrap\_example}
%
\begin{Description}
Output example for a bootstrapping analysis conducted with the ResIN package
\end{Description}
%
\begin{Usage}
\begin{verbatim}
data(Bootstrap_example)
\end{verbatim}
\end{Usage}
%
\begin{Format}
An object of class \code{"rda"}
\end{Format}
%
\begin{References}
This dataset was made available by Lüders et.al. 2024.
\end{References}
%
\begin{Examples}
\begin{ExampleCode}

data(Bootstrap_example)
str(Bootstrap_example)

\end{ExampleCode}
\end{Examples}
\HeaderA{BrJSocPsychol\_2024}{Source data for Lüders, A., Carpentras, D. and Quayle, M., 2024. Attitude networks as intergroup realities: Using network‐modelling to research attitude‐identity relationships in polarized political contexts. British Journal of Social Psychology, 63(1), pp.37-51.}{BrJSocPsychol.Rul.2024}
\keyword{datasets}{BrJSocPsychol\_2024}
%
\begin{Description}
The sample of N = 402 paid participants through the crowd working platform Prolific Academic. The core of the survey consists of A set of eight political attitude items abortion, immigration, gun control, and gay marriage. Each item followed a 5-point scale ranging from strong disagreement to strong agreement. The survey also includes items on partisanship, affective polarization, and a short vignette experiment.
\end{Description}
%
\begin{Usage}
\begin{verbatim}
data(BrJSocPsychol_2024)
\end{verbatim}
\end{Usage}
%
\begin{Format}
An object of class \code{"data.frame"}
\end{Format}
%
\begin{References}
This dataset was made available by Lüders et.al. 2024.
\end{References}
%
\begin{Examples}
\begin{ExampleCode}

data(BrJSocPsychol_2024)
head(BrJSocPsychol_2024)


\end{ExampleCode}
\end{Examples}
\HeaderA{lik\_data}{Likert-type, mock-response data for "ResIN" package examples}{lik.Rul.data}
\keyword{datasets}{lik\_data}
%
\begin{Description}
An artificially created data-set (n=1000) of 12, 5-point Likert data. Modeled on the basis of a standard normal data-generating process. Likert scales contain 20 percent uncorrelated, homoscedastic measurement error. This data-set is used for the examples in the "ResIN" package vignette.
\end{Description}
%
\begin{Usage}
\begin{verbatim}
data(lik_data)
\end{verbatim}
\end{Usage}
%
\begin{Format}
An object of class \code{"data.frame"}
\end{Format}
%
\begin{References}
This data set was artificially created for the ResIN package.
\end{References}
%
\begin{Examples}
\begin{ExampleCode}

data(lik_data)
head(lik_data)


\end{ExampleCode}
\end{Examples}
\HeaderA{plot.ResIN}{Plot a ResIN object}{plot.ResIN}
%
\begin{Description}
Plot a ResIN object
\end{Description}
%
\begin{Usage}
\begin{verbatim}
## S3 method for class 'ResIN'
plot(x, which = c("network", "bipartite"), print_plot = TRUE, ...)
\end{verbatim}
\end{Usage}
%
\begin{Arguments}
\begin{ldescription}
\item[\code{x}] A ResIN object.

\item[\code{which}] Which plot to show: \code{"network"} (default) or \code{"bipartite"}.

\item[\code{print\_plot}] Logical; print the plot (TRUE) or just return it (FALSE).

\item[\code{...}] Ignored.
\end{ldescription}
\end{Arguments}
%
\begin{Value}
The plot object (invisibly if \code{print\_plot = TRUE}).
\end{Value}
\HeaderA{print.ResIN}{Print a ResIN object}{print.ResIN}
%
\begin{Description}
Print a ResIN object
\end{Description}
%
\begin{Usage}
\begin{verbatim}
## S3 method for class 'ResIN'
print(x, ...)
\end{verbatim}
\end{Usage}
%
\begin{Arguments}
\begin{ldescription}
\item[\code{x}] A ResIN object.

\item[\code{...}] Ignored.
\end{ldescription}
\end{Arguments}
\HeaderA{ResIN}{Flagship function that implements Response Item Network (ResIN) analysis}{ResIN}
%
\begin{Description}
Performs Response Item-Network (ResIN) analysis in one go. Users minimally need to supply a dataframe or matrix of discrete response data. If needed for step-wise analysis, all intermediate outputs can still be accessed as part of the aux\_objects output list.
\end{Description}
%
\begin{Usage}
\begin{verbatim}
ResIN(
  df,
  node_vars = NULL,
  left_anchor = NULL,
  cor_method = "pearson",
  weights = NULL,
  missing_cor = "pairwise",
  offset = 0,
  ResIN_scores = TRUE,
  remove_nonsignificant = FALSE,
  remove_nonsignificant_method = "default",
  sign_threshold = 0.05,
  node_covars = NULL,
  node_costats = NULL,
  network_stats = TRUE,
  detect_clusters = FALSE,
  cluster_method = NULL,
  cluster_arglist = NULL,
  cluster_assignment = TRUE,
  generate_ggplot = TRUE,
  plot_ggplot = TRUE,
  plot_whichstat = NULL,
  plot_edgestat = NULL,
  color_palette = "RdBu",
  direction = 1,
  plot_responselabels = TRUE,
  response_levels = NULL,
  plot_title = NULL,
  bipartite = FALSE,
  save_input = TRUE,
  remove_negative = TRUE,
  EBICglasso = FALSE,
  EBICglasso_arglist = NULL,
  seed = NULL
)
\end{verbatim}
\end{Usage}
%
\begin{Arguments}
\begin{ldescription}
\item[\code{df}] A data-frame object containing the raw data.

\item[\code{node\_vars}] An optional character vector detailing the attitude item columns to be selected for ResIN analysis (i.e. the subset of attitude variables in df).

\item[\code{left\_anchor}] An optional character scalar indicating a particular response node which determines the spatial orientation of the ResIN latent space. If this response node does not appear on the left-hand side, the x-plane will be inverted. This ensures consistent interpretation of the latent space across multiple iterations (e.g. in bootstrapping analysis). Defaults to NULL (no adjustment to orientation is taken.)

\item[\code{cor\_method}] Which correlation method should be used? Current implementation supports "pearson" (default) and "polychoric". Please note that polychoric correlations are currently unsupported for weighted analysis.

\item[\code{weights}] Optional survey weights. Can be either \code{NULL} (default), a numeric vector of length \code{nrow(df)}, or a character scalar naming a weights column in \code{df}. If a column name is supplied and \code{node\_vars = NULL}, the weights column is automatically excluded from the response-node variables used for ResIN estimation.

\item[\code{missing\_cor}] Character scalar controlling missing-data handling for correlation estimation. Either \code{"pairwise"} (default) or  \code{"listwise"}.

\item[\code{offset}] Optional off-set to correlation edges to manually adjust for over- or under-fitting the network. Defaults to \code{0}. Supplying a value between -1 and 0 globally reduces edge values by that amount, leading to the elimination of all positive  edges below that value, resulting in a more sparse network. (However, we strongly recommend setting \code{remove\_nonsignificant=TRUE} instead for a more principled approach to ensuring optimal network sparsity as global thresholds have heuristic value at best). Alternatively, a value between 0 and 1 enforces a positive offset, resulting in more dense (but potentially over-fitted) networks.

\item[\code{ResIN\_scores}] Logical; should spatial scores be calculated for every individual. Defaults to TRUE. Function obtains the mean positional score on the major (x-axis) and minor (y-axis). Current package implementation also provides empirical Bayesian scores via James-Stein shrinkage (\code{eb\_x}) and heuristic shrinkage (\code{heur\_x}) scores. Please refer to the package [vignette]https://pwarncke77.github.io/ResIN/articles/ResIN-VIGNETTE.html\#spatial-interpretation-and-individual-latent-space-scores for further details.

\item[\code{remove\_nonsignificant}] Logical; should non-significant edges be removed from the ResIN network? Defaults to FALSE. For weighted Pearson correlations, p-values are approximated using a weighted effective sample size. For currently unsupported polychoric configurations, ResIN falls back to Pearson and issues a warning.

\item[\code{remove\_nonsignificant\_method}] Character scalar specifying how p-values are thresholded when \code{remove\_nonsignificant = TRUE}. Defaults to \code{"default"}, which prunes edges with raw p-values greater than \code{sign\_threshold} (i.e., retains edges with \code{p <= sign\_threshold}). If set to \code{"bh"}, p-values are adjusted using the Benjamini--Hochberg procedure; edges are retained only if the adjusted p-value is less than or equal to \code{sign\_threshold}, interpreted as the target false discovery rate \eqn{q}{}. This provides multiplicity control across all tested edges and is typically more principled than using unadjusted p-values, but may be slightly slower. See \code{\LinkA{p.adjust}{p.adjust}} for details.

\item[\code{sign\_threshold}] Numeric scalar controlling the pruning threshold used when \code{remove\_nonsignificant = TRUE}. For \code{remove\_nonsignificant\_method = "default"}, this is the raw p-value cutoff (e.g., \code{0.05}). For \code{remove\_nonsignificant\_method = "bh"}, this is the target false discovery rate level \eqn{q}{} (e.g., \code{0.05}), applied to Benjamini--Hochberg adjusted p-values.

\item[\code{node\_covars}] An optional character string selecting quantitative co-variates that can be used to enhance ResIN analysis. Typically, these covariates provide grouped summary statistics for item response nodes. (E.g.: What is the average age or income level of respondents who selected a particular item response?) Variable names specified here should match existing columns in \code{df}.

\item[\code{node\_costats}] If any \code{node\_covars} are selected, what summary statistics should be estimated from them? Argument should be a character vector and call a base-R function. (E.g. \code{"mean"}, \code{"median"}, \code{"sd"}). Each element specified in \code{node\_costats} is applied to each element in \code{node\_covars} and the out-put is stored as a node-level summary statistic in the \code{ResIN\_nodeframe}. The extra columns in \code{ResIN\_nodeframe} are labeled according to the following template: "covariate name"\_"statistic". So for the respondents mean age, the corresponding column in \code{ResIN\_nodeframe} would be labeled as "age\_mean".

\item[\code{network\_stats}] Should common node- and graph level network statistics be extracted? Calls \code{qgraph::centrality\_auto} and \code{DirectedClustering::ClustF} to the ResIN graph object to extract node-level betweenness, closeness, strength centrality, as well as the mean and standard deviation of these scores at the network level. Also estimates network expected influence, average path length, and global clustering coefficients. Defaults to TRUE. Set to FALSE if estimation takes a long time.

\item[\code{detect\_clusters}] Optional, should community detection be performed on item response network? Defaults to FALSE. If set to TRUE, performs a clustering method from the [igraph](https://igraph.org/r/doc/cluster\_leading\_eigen.html) library and stores the results in the \code{ResIN\_nodeframe} output.

\item[\code{cluster\_method}] A character scalar specifying the [igraph-based](https://igraph.org/r/doc/communities.html) community detection function.

\item[\code{cluster\_arglist}] An optional list specifying additional arguments to the selected [igraph](https://igraph.org/r/doc/communities.html) clustering method.

\item[\code{cluster\_assignment}] Should individual (survey) respondents be assigned to different clusters? If set to TRUE, function will generate an n*c matrix of probabilities for each respondent to be assigned to one of c clusters. Furthermore, a vector of length n is generated displaying the most likely cluster respondents belong to. In case of a tie between one or more clusters, a very small amount of random noise determines assignment. Both matrix and vectors are added to the \code{aux\_objects} list. Defaults to FALSE and will be ignored if \code{detect\_clusters} is set to FALSE.

\item[\code{generate\_ggplot}] Logical; should a ggplot-based visualization of the ResIN network be generated? Defaults to TRUE.

\item[\code{plot\_ggplot}] Logical; should a basic ggplot of the ResIN network be plotted? Defaults to TRUE. If set to FALSE, the ggplot object will not be directly returned to the console. (However, if generate\_ggplot=TRUE, the plot will still be generated and stored alongside the other output objects.)

\item[\code{plot\_whichstat}] Should a particular node-level metric be color-visualized in the ggplot output? For node cluster, specify "cluster". For the same Likert response choices or options, specify "choices". For a particular node-level co-variate please specify the name of the particular element in \code{node\_covars} followed by a "\_" and the specific \code{node\_costats} you would like to visualize. For instance if you want the visualize average age at the node-level, you should specify "age\_mean". To colorize by node centrality statistics, possible choices are "Strength", "Betweenness", "Closeness", and "ExpectedInfluence". Defaults to NULL. Make sure to supply appropriate choices to \code{node\_covars}, \code{node\_costats}, \code{detect\_clusters}, and/or \code{network\_stats} prior to setting this argument.

\item[\code{plot\_edgestat}] Should the thickness of the edges be adjusted according to a particular co-statistic? Defaults to NULL. Possible choices are "weight" for the bi-variate correlation strength, and "edgebetweenness"

\item[\code{color\_palette}] Optionally, you may specify the ggplot2 color palette to be applied to the plot. All options contained in [\code{RColorBrewer}](https://cran.r-project.org/web/packages/RColorBrewer/RColorBrewer.pdf) (for discrete colors such as cluster assignments) and [\code{ggplot2::scale\_colour\_distiller}](https://ggplot2.tidyverse.org/reference/scale\_brewer.html) are supported. Defaults to "RdBu".

\item[\code{direction}] Which direction should the color palette be applied in? Defaults to 1. Set to -1 if the palette should appear in reverse order.

\item[\code{plot\_responselabels}] Should response labels be plotted via \code{geom\_text}? Defaults to TRUE. It is recommended to set to FALSE if the network possesses a lot of nodes and/or long response choice names.

\item[\code{response\_levels}] An optional character vector specifying the correct order of global response levels. Only useful if all node-items follow the same convention (e.g. ranging from "strong disagreement" to "strong agreement"). The supplied vector should have the same length as the total number of response options and supply these (matching exactly) in the correct order. E.g. c("Strongly Agree", "Somewhat Agree", "Neutral", "Somewhat Disagree", "Strongly Disagree"). Defaults to NULL.

\item[\code{plot\_title}] Optionally, a character scalar specifying the title of the ggplot output. Defaults to "ResIN plot".

\item[\code{bipartite}] Logical; should a bipartite graph be produced in addition to classic ResIN graph? Defaults to FALSE. If set to TRUE, an  [igraph](https://igraph.org/r/doc/) bipartite graph with response options as node type 1 and participants as node type 2 will be generated and included in the output list. Further, an object called \code{coordinate\_df} with spatial coordinates of respondents and a plot-able \code{ggraph}-object called \code{bipartite\_ggraph} are generated if set to TRUE.

\item[\code{save\_input}] Logical; should input data and function arguments be saved (this is necessary for running ResIN\_boots\_prepare function). Defaults to TRUE.

\item[\code{remove\_negative}] Logical; should all negative correlations be removed? Defaults to TRUE (highly recommended). Setting to FALSE makes it impossible to estimate a force-directed network layout. Function will use igraph::layout\_nicely instead.

\item[\code{EBICglasso}] Retired as of ResIN 2.3.0 and ignored.

\item[\code{EBICglasso\_arglist}] Retired as of ResIN 2.3.0 and ignored.

\item[\code{seed}] Random seed for force-directed algorithm. Defaults to NULL (no seed is set.) If scalar integer is supplied, that seed will be set prior to analysis.
\end{ldescription}
\end{Arguments}
%
\begin{Value}
An edge-list type data-frame, \code{ResIN\_edgelist}, a node-level data-frame, \code{ResIN\_nodeframe}, an n*2 data-frame of individual-level spatial scores along the major (x) and minor(y) axis, \code{ResIN\_scores} a list of graph-level statistics \code{graph\_stats} including (\code{graph\_structuration}), and centralization (\code{graph\_centralization}). Further, a \code{bipartite\_output} list which includes an \code{igraph} class bipartite graph (\code{bipartite\_igraph}), a data frame, \code{coordinate\_df}, with spatial coordinates of respondents, and a plot-able \code{ggraph}-object called \code{bipartite\_ggraph} is optionally generated. Lastly, the output includes a list of auxiliary objects, \code{aux\_objects}, including the ResIN adjacency matrix (\code{adj\_matrix}), a numeric vector detailing which item responses belong to which item (\code{same\_items}), and the dummy-coded item-response data-frame (\code{df\_dummies}). For reproducibility, (\code{aux\_objects\$meta} stores a numeric dataframe identifier (\code{df\_id}, the random seed, call, and the (\code{ResIN} package version used to create the object.”
\end{Value}
%
\begin{Examples}
\begin{ExampleCode}

## Load the 12-item simulated Likert-type toy dataset
data(lik_data)

# Apply the ResIN function to toy Likert data:
ResIN_obj <- ResIN(lik_data, network_stats = TRUE, remove_nonsignificant = TRUE)

\end{ExampleCode}
\end{Examples}
\HeaderA{ResIN-coercion-generics}{Internal coercion generics}{ResIN.Rdash.coercion.Rdash.generics}
\aliasA{as.gephi}{ResIN-coercion-generics}{as.gephi}
\aliasA{as.qgraph}{ResIN-coercion-generics}{as.qgraph}
\aliasA{as.tidygraph}{ResIN-coercion-generics}{as.tidygraph}
\keyword{internal}{ResIN-coercion-generics}
%
\begin{Description}
These functions are S3 generics used for object coercion in \pkg{ResIN}.
End users should typically consult the class methods:
\code{\LinkA{as.igraph.ResIN}{as.igraph.ResIN}}, \code{\LinkA{as.qgraph.ResIN}{as.qgraph.ResIN}}, and \code{\LinkA{as.gephi.ResIN}{as.gephi.ResIN}}.
\end{Description}
%
\begin{Usage}
\begin{verbatim}
as.qgraph(x, ...)

as.gephi(x, ...)

as.tidygraph(x, ...)
\end{verbatim}
\end{Usage}
\HeaderA{ResIN-reexports}{Re-exported functions used by ResIN}{ResIN.Rdash.reexports}
\aliasA{as.igraph}{ResIN-reexports}{as.igraph}
\aliasA{as.network}{ResIN-reexports}{as.network}
\keyword{internal}{ResIN-reexports}
%
\begin{Description}
These functions are re-exported so users can call common coercion generics
(e.g., \code{as.igraph()} and \code{as.network()}) directly after loading
\pkg{ResIN}, with S3 dispatch to \code{ResIN} methods.
\end{Description}
%
\begin{Usage}
\begin{verbatim}
as.igraph(x, ...)

as.network(x, ...)
\end{verbatim}
\end{Usage}
\HeaderA{ResIN\_boots\_draws}{Extract bootstrap draws from ResIN objects}{ResIN.Rul.boots.Rul.draws}
%
\begin{Description}
\code{"ResIN\_boots\_draws"} objects are typically returned by \code{\LinkA{ResIN\_boots\_extract}{ResIN.Rul.boots.Rul.extract}}
when the requested quantity is scalar per bootstrap iteration (e.g., \code{"global\_clustering"}).
The object is a numeric vector with attributes describing the bootstrap context.
\end{Description}
%
\begin{Details}
Common attributes include:
\begin{description}

\item[\code{what}] Name of the extracted quantity.
\item[\code{n\_total}, \code{n\_ok}, \code{n\_failed}] Counts of total, successful, and failed iterations.
\item[\code{plan}] The \code{"ResIN\_boots\_prepped"} plan used to generate the results (if attached).
\item[\code{created}] Timestamp of execution (if attached).

\end{description}

\end{Details}
%
\begin{Section}{Methods}

\begin{description}

\item[\code{print(x)}] Compact display including median and 95\% CI.
\item[\code{summary(object)}] Return descriptive statistics and quantiles.
\item[\code{confint(object)}] Quantile-based confidence intervals.
\item[\code{plot(x)}] Histogram with CI markers.

\end{description}

\end{Section}
\HeaderA{ResIN\_boots\_execute}{Carry out prepared bootstrap analyses on ResIN networks}{ResIN.Rul.boots.Rul.execute}
%
\begin{Description}
Executes a bootstrap plan created by \code{\LinkA{ResIN\_boots\_prepare}{ResIN.Rul.boots.Rul.prepare}} by repeatedly
re-estimating ResIN on resampled or permuted versions of the original data.
Can optionally leverage CPU parallelism.
\end{Description}
%
\begin{Usage}
\begin{verbatim}
ResIN_boots_execute(
  ResIN_boots_prepped,
  parallel = FALSE,
  detect_cores = TRUE,
  core_offset = 0L,
  n_cores = 2L,
  inorder = FALSE
)
\end{verbatim}
\end{Usage}
%
\begin{Arguments}
\begin{ldescription}
\item[\code{ResIN\_boots\_prepped}] A \code{"ResIN\_boots\_prepped"} bootstrap plan (output of \code{\LinkA{ResIN\_boots\_prepare}{ResIN.Rul.boots.Rul.prepare}}).

\item[\code{parallel}] Should execution use parallelism via \code{foreach} + a PSOCK cluster? Defaults to FALSE.

\item[\code{detect\_cores}] Should available CPU cores be detected automatically? Defaults to TRUE (ignored when \code{parallel = FALSE}).

\item[\code{core\_offset}] Integer offset subtracted from the number of detected cores. Defaults to 0L.

\item[\code{n\_cores}] Manually specify number of cores (ignored if \code{detect\_cores = TRUE} or \code{parallel = FALSE}).

\item[\code{inorder}] Should parallel execution preserve sequential ordering? Defaults to FALSE.
\end{ldescription}
\end{Arguments}
%
\begin{Value}
An object of class \code{"ResIN\_boots\_executed"} containing \code{n} bootstrapped
\code{ResIN} fits. Use \code{print()}, \code{summary()}, \code{length()}, and \code{[}
to inspect or subset results. See \code{\LinkA{ResIN\_boots\_executed}{ResIN.Rul.boots.Rul.executed}} for details.
\end{Value}
%
\begin{Examples}
\begin{ExampleCode}
## Load the 12-item simulated Likert-type toy dataset
data(lik_data)

# Apply the ResIN function to toy Likert data:
ResIN_obj <- ResIN(lik_data, network_stats = TRUE,
                      generate_ggplot = FALSE, plot_ggplot = FALSE)


# Prepare for bootstrapping
prepped_boots <- ResIN_boots_prepare(ResIN_obj, n=100, boots_type="resample")

# Execute the prepared bootstrap list
executed_boots <-  ResIN_boots_execute(prepped_boots, parallel = TRUE, detect_cores = TRUE)

# Extract results - here for example, the network (global)-clustering coefficient
ResIN_boots_extract(executed_boots, what = "global_clustering", summarize_results = TRUE)


\end{ExampleCode}
\end{Examples}
\HeaderA{ResIN\_boots\_executed}{Bootstrapped results}{ResIN.Rul.boots.Rul.executed}
%
\begin{Description}
The \code{"ResIN\_boots\_executed"} object is returned by
\code{\LinkA{ResIN\_boots\_execute}{ResIN.Rul.boots.Rul.execute}}. It contains a list of fitted \code{"ResIN"} objects,
typically produced by refitting ResIN on resampled/permuted versions of the original data.
\end{Description}
%
\begin{Details}
The object is a list of length \code{n} where each element is (typically) a
\code{\LinkA{ResIN}{ResIN}} fit. For reproducibility, the following attributes may be present:
\begin{description}

\item[\code{plan}] A \code{"ResIN\_boots\_prepped"} plan used to generate the results.
\item[\code{boot\_inputs}] Optional list of bootstrap input data sets (only if \code{save\_input = TRUE} in the plan).
\item[\code{created}] Time stamp when the executed object was created.

\end{description}

\end{Details}
%
\begin{Section}{Methods}

\begin{description}

\item[\code{print(x)}] Print a compact overview of the bootstrap results.
\item[\code{summary(object)}] Summarize the results (iterations, successes/failures, plan details).
\item[\code{length(x)}] Return the number of iterations.
\item[\code{x[i]}] Subset the results while preserving attached attributes.

\end{description}

\end{Section}
%
\begin{SeeAlso}
\code{\LinkA{ResIN\_boots\_prepare}{ResIN.Rul.boots.Rul.prepare}}, \code{\LinkA{ResIN\_boots\_execute}{ResIN.Rul.boots.Rul.execute}},
\code{\LinkA{ResIN\_boots\_extract}{ResIN.Rul.boots.Rul.extract}}, \code{\LinkA{ResIN\_boots\_prepped}{ResIN.Rul.boots.Rul.prepped}}
\end{SeeAlso}
\HeaderA{ResIN\_boots\_extract}{Extract and summarize the results of a bootstrap simulation undertaken on a ResIN object}{ResIN.Rul.boots.Rul.extract}
%
\begin{Description}
Extract bootstrap draws from \code{"ResIN\_boots\_executed"} results. Failed iterations
(stored as \code{NULL}) are skipped automatically.
\end{Description}
%
\begin{Usage}
\begin{verbatim}
ResIN_boots_extract(
  ResIN_boots_executed,
  what,
  summarize_results = FALSE,
  allow_missing = FALSE
)
\end{verbatim}
\end{Usage}
%
\begin{Arguments}
\begin{ldescription}
\item[\code{ResIN\_boots\_executed}] An object of class \code{"ResIN\_boots\_executed"} (output of \code{\LinkA{ResIN\_boots\_execute}{ResIN.Rul.boots.Rul.execute}}).

\item[\code{what}] Character scalar naming the quantity to extract (e.g., \code{"global\_clustering"}).

\item[\code{summarize\_results}] Logical; if TRUE return a small summary table, otherwise return draws.

\item[\code{allow\_missing}] Logical; if FALSE (default) the function errors if \code{what} is missing in any successful fit.
If TRUE, missing iterations are kept as \code{NULL} and summaries are computed from available values.
\end{ldescription}
\end{Arguments}
%
\begin{Value}
If the extracted quantity is scalar per iteration, returns an object of class
\code{"ResIN\_boots\_draws"} (a numeric vector with attributes). If \code{summarize\_results = TRUE},
returns a one-row data.frame of summary statistics.
\end{Value}
%
\begin{Examples}
\begin{ExampleCode}
## Load the 12-item simulated Likert-type toy dataset
data(lik_data)

# Apply the ResIN function to toy Likert data:
ResIN_obj <- ResIN(lik_data, network_stats = TRUE,
                      generate_ggplot = FALSE, plot_ggplot = FALSE)


# Prepare for bootstrapping
prepped_boots <- ResIN_boots_prepare(ResIN_obj, n=100, boots_type="resample")

# Execute the prepared bootstrap list
executed_boots <-  ResIN_boots_execute(prepped_boots, parallel = TRUE, detect_cores = TRUE)

# Extract results - here for example, the network (global)-clustering coefficient
ResIN_boots_extract(executed_boots, what = "global_clustering", summarize_results = TRUE)



\end{ExampleCode}
\end{Examples}
\HeaderA{ResIN\_boots\_prepare}{Create a bootstrap plan for re-estimating ResIN objects to derive statistical uncertainty estimates}{ResIN.Rul.boots.Rul.prepare}
%
\begin{Description}
Provides instructions for how to bootstrap a ResIN network to derive uncertainty estimates around core quantities of interest. Requires output of \code{ResIN} function.
\end{Description}
%
\begin{Usage}
\begin{verbatim}
ResIN_boots_prepare(
  ResIN_object,
  n = 10000,
  boots_type = "resample",
  resample_size = NULL,
  weights = NULL,
  save_input = FALSE,
  seed_boots = 42
)
\end{verbatim}
\end{Usage}
%
\begin{Arguments}
\begin{ldescription}
\item[\code{ResIN\_object}] A ResIN object to prepare bootstrapping workflow.

\item[\code{n}] Bootstrapping sample size. Defaults to 10.000.

\item[\code{boots\_type}] What kind of bootstrapping should be performed? If set to "resample", function performs row-wise re-sampling of raw data (useful for e.g., sensitivity or power analysis). If set to "permute", function will randomly reshuffle raw item responses (useful e.g., for simulating null-hypothesis distributions). Defaults to "resample".

\item[\code{resample\_size}] Optional parameter determining sample size when \code{boots\_type} is set to "resample". Defaults of to number of rows in raw data.

\item[\code{weights}] An optional weights vector that can be used to adjust the re-sampling of observations. Should either be NULL (default) or a positive numeric vector of the same length as the original data.

\item[\code{save\_input}] Should all input information for each bootstrap iteration (including re-sampled/permuted data) be stored. Set to FALSE by default to save a lot of memory and disk storage.

\item[\code{seed\_boots}] Random seed for bootstrap samples
\end{ldescription}
\end{Arguments}
%
\begin{Value}
An object of class \code{"ResIN\_boots\_prepped"} containing a bootstrap plan
(specification) used by \code{\LinkA{ResIN\_boots\_execute}{ResIN.Rul.boots.Rul.execute}}.
Use \code{print()}, \code{summary()}, \code{length()}, and \code{[}
to inspect or subset the plan. See \code{\LinkA{ResIN\_boots\_prepped}{ResIN.Rul.boots.Rul.prepped}} for details.
\end{Value}
%
\begin{Examples}
\begin{ExampleCode}
## Load the 12-item simulated Likert-type toy dataset
data(lik_data)

# Apply the ResIN function to toy Likert data:
ResIN_obj <- ResIN(lik_data, network_stats = TRUE,
                      generate_ggplot = FALSE, plot_ggplot = FALSE)


# Prepare for bootstrapping
prepped_boots <- ResIN_boots_prepare(ResIN_obj, n=100, boots_type="resample")

# Execute the prepared bootstrap list
executed_boots <-  ResIN_boots_execute(prepped_boots, parallel = TRUE, detect_cores = TRUE)

# Extract results - here for example, the network (global)-clustering coefficient
ResIN_boots_extract(executed_boots, what = "global_clustering", summarize_results = TRUE)


\end{ExampleCode}
\end{Examples}
\HeaderA{ResIN\_boots\_prepped}{Prepared bootstrap plan}{ResIN.Rul.boots.Rul.prepped}
%
\begin{Description}
The \code{"ResIN\_boots\_prepped"} object is created by
\code{\LinkA{ResIN\_boots\_prepare}{ResIN.Rul.boots.Rul.prepare}} and provides a reproducible specification
for running a bootstrap with \code{\LinkA{ResIN\_boots\_execute}{ResIN.Rul.boots.Rul.execute}}.
\end{Description}
%
\begin{Details}
A \code{"ResIN\_boots\_prepped"} object is a list with (at least) the following elements:
\begin{description}

\item[call] The matched call used to create the plan.
\item[boots\_type] Character; \code{"resample"} or \code{"permute"}.
\item[n] Integer; number of bootstrap iterations.
\item[resample\_size] Integer; sample size used when resampling rows.
\item[weights] Optional numeric vector of sampling weights (or \code{NULL}).
\item[save\_input] Logical; whether to store bootstrap inputs during execution.
\item[seed\_boots] Integer; seed used to generate per-iteration seeds.
\item[iter\_seeds] Integer vector; per-iteration RNG seeds (useful for parallel execution).
\item[arglist] A list of arguments passed to \code{\LinkA{ResIN}{ResIN}} for each re-fit
(typically based on the original ResIN call, with plotting disabled).
\item[df\_id] Character; MD5 hash identifying the raw data used to create the plan.
\item[ResIN\_version] Character; ResIN package version used to create the plan.

\end{description}

\end{Details}
%
\begin{Section}{Methods}

\begin{description}

\item[\code{print(x)}] Print a compact summary of the bootstrap plan.
\item[\code{summary(object)}] Return a structured summary of the plan.
\item[\code{length(x)}] Return the number of iterations \code{n}.
\item[\code{x[i]}] Subset the plan to selected iteration indices (also updates \code{n}).

\end{description}

\end{Section}
%
\begin{SeeAlso}
\code{\LinkA{ResIN\_boots\_prepare}{ResIN.Rul.boots.Rul.prepare}}, \code{\LinkA{ResIN\_boots\_execute}{ResIN.Rul.boots.Rul.execute}},
\code{\LinkA{ResIN\_boots\_extract}{ResIN.Rul.boots.Rul.extract}}
\end{SeeAlso}
\HeaderA{ResIN\_to\_gephi}{(Deprecated.) Convert ResIN networks to Gephi-readable csv tables. Use \code{as.gephi()} method instead.}{ResIN.Rul.to.Rul.gephi}
%
\begin{Description}
Deprecated/legacy function. Saves a ResIN graph as a series of csv files readable by Gephi. Now supplanted by \code{as.gephi()} method.
\end{Description}
%
\begin{Usage}
\begin{verbatim}
ResIN_to_gephi(
  ResIN_object,
  file = "ResIN_gephi.csv",
  edges_only = TRUE,
  dont_save_csv = FALSE
)
\end{verbatim}
\end{Usage}
%
\begin{Arguments}
\begin{ldescription}
\item[\code{ResIN\_object}] The output of the ResIN function (a list with class ResIN).

\item[\code{file}] The name with .csv extension for the Gephi readable file to be output at. Defaults to "ResIN\_gephi.csv".

\item[\code{edges\_only}] Logical; if TRUE write/return only edges.

\item[\code{dont\_save\_csv}] Logical; set TRUE to disable writing.
\end{ldescription}
\end{Arguments}
%
\begin{Value}
A series of csv files readable by Gephi
\end{Value}
%
\begin{References}
Source code of original function (< version 2.2.0) had been adapted from: https://github.com/RMHogervorst/gephi?tab=MIT-1-ov-file\#readme
\end{References}
%
\begin{Examples}
\begin{ExampleCode}
## Load the 12-item simulated Likert-type ResIN toy dataset
data(lik_data)

## Estimate a ResIN network
res <- ResIN(lik_data, generate_ggplot = FALSE)

## Create Gephi edge table without writing files
edges <- as.gephi(res, dont_save_csv = TRUE)
head(edges)

## Not run: 
## Write CSV file(s) for import to Gephi
## (writes "ResIN_gephi.csv" by default)
as.gephi(res, file = "ResIN_gephi.csv")

## Write both edges and nodes tables
## (writes "ResIN_gephi_edges.csv" and "ResIN_gephi_nodes.csv")
as.gephi(res, file = "ResIN_gephi.csv", edges_only = FALSE)

## End(Not run)


\end{ExampleCode}
\end{Examples}
\HeaderA{ResIN\_to\_igraph}{(Deprecated.) Convert a ResIN network into an igraph object. Use \code{as.igraph()} method instead.}{ResIN.Rul.to.Rul.igraph}
%
\begin{Description}
Deprecated/legacy function. Transforms the output of the \code{ResIN} function into an [igraph](https://igraph.org/r/doc/cluster\_leading\_eigen.html) object. Now simply a wrapper for the \code{as.igraph()} method.
\end{Description}
%
\begin{Usage}
\begin{verbatim}
ResIN_to_igraph(ResIN_object, igraph_arglist = NULL)
\end{verbatim}
\end{Usage}
%
\begin{Arguments}
\begin{ldescription}
\item[\code{ResIN\_object}] the output of the ResIN function (a list with class ResIN).

\item[\code{igraph\_arglist}] an optional argument list to be supplied to the igraph::graph\_from\_adjacency\_matrix function. If NULL, default is: list(mode = "undirected", weighted = TRUE, diag = FALSE).
\end{ldescription}
\end{Arguments}
%
\begin{Value}
A class \code{igraph} object.
\end{Value}
%
\begin{References}
Csardi G, Nepusz T (2006). “The igraph software package for complex network research.” InterJournal, Complex Systems, 1695. https://igraph.org.
\end{References}
%
\begin{SeeAlso}
\code{\LinkA{as.igraph}{as.igraph.ResIN}} as the recommended interface.
\end{SeeAlso}
%
\begin{Examples}
\begin{ExampleCode}

## Load the 12-item simulated Likert-type ResIN toy dataset
data(lik_data)

## Run the function:

igraph_output <-  as.igraph(ResIN(lik_data, plot_ggplot = FALSE))

class(igraph_output)

## Plot and/or investigate as you wish:

igraph::plot.igraph(igraph_output)



\end{ExampleCode}
\end{Examples}
\HeaderA{ResIN\_to\_qgraph}{(Deprecated.) Convert a ResIN network into an qgraph object. Use \code{as.qgraph()} method instead.}{ResIN.Rul.to.Rul.qgraph}
%
\begin{Description}
Deprecated/legacy function. Transforms the output of the \code{ResIN} function into an \code{qgraph} object. Use \code{as.qgraph()} method instead.
\end{Description}
%
\begin{Usage}
\begin{verbatim}
ResIN_to_qgraph(ResIN_object, qgraph_arglist = NULL)
\end{verbatim}
\end{Usage}
%
\begin{Arguments}
\begin{ldescription}
\item[\code{ResIN\_object}] the output of the ResIN function (a list with class ResIN).

\item[\code{qgraph\_arglist}] an optional argument list to be supplied to the igraph::graph\_from\_adjacency\_matrix function. If NULL, defaults are: \code{list(layout = "spring", maximum = 1, vsize = 6, DoNotPlot = TRUE, sampleSize = nrow(df\_nodes), mar = c(3,3,3,3), normalize = FALSE)}
\end{ldescription}
\end{Arguments}
%
\begin{Value}
A [qgraph]https://cran.r-project.org/web/packages/qgraph/index.html graph object.
\end{Value}
%
\begin{References}
Epskamp S, Cramer AOJ, Waldorp LJ, Schmittmann VD, Borsboom D (2012). “qgraph: Network Visualizations of Relationships in Psychometric Data.” Journal of Statistical Software, 48(4), 1–18.
\end{References}
%
\begin{SeeAlso}
\code{\LinkA{as.qgraph}{as.qgraph.ResIN}} as the recommended interface.
\end{SeeAlso}
%
\begin{Examples}
\begin{ExampleCode}
## Load the 12-item simulated Likert-type ResIN toy dataset
data(lik_data)

## Run the function:
ResIN_qgraph <-  as.qgraph(ResIN(lik_data, plot_ggplot = FALSE))

class(ResIN_qgraph)

\end{ExampleCode}
\end{Examples}
\HeaderA{summary.ResIN}{Summarize a ResIN object}{summary.ResIN}
%
\begin{Description}
Summarize a ResIN object
\end{Description}
%
\begin{Usage}
\begin{verbatim}
## S3 method for class 'ResIN'
summary(object, ...)
\end{verbatim}
\end{Usage}
%
\begin{Arguments}
\begin{ldescription}
\item[\code{object}] A ResIN object.

\item[\code{...}] Ignored.
\end{ldescription}
\end{Arguments}
%
\begin{Value}
An object of class \code{summary.ResIN}.
\end{Value}
\printindex{}
\end{document}
